%! AP Statistics — Unit 9, Lesson 9.4 — Student Packet Cover Sheet
\documentclass[10pt]{article}
\usepackage{apstats-article}
\usepackage{apstats-boxes}
\usepackage{tikz}

\begin{document}

\begin{tikzpicture}[remember picture, overlay]
  \fill[navy] (current page.north west) rectangle
    ([yshift=-2.2cm]current page.north east);
  \node[anchor=west, white, font=\large\bfseries] at
    ([xshift=1cm, yshift=-0.7cm]current page.north west)
    {\CourseName: \SchoolYear};
  \node[anchor=west, white, font=\normalsize] at
    ([xshift=1cm, yshift=-1.3cm]current page.north west)
    {Unit 9: Inference for Quantitative Data: Slopes};
  \node[anchor=west, white, font=\normalsize\bfseries] at
    ([xshift=1cm, yshift=-1.9cm]current page.north west)
    {Lesson 9.4 \quad Setting Up a Test for the Slope of a Regression Model};
\end{tikzpicture}

\vspace{2.4cm}
\namedateperiod

\begin{learningtargetbox}
\begin{itemize}
    \item I can identify a $t$-test for slope as the appropriate procedure when
          testing a claim about $\beta$.
    \item I can state H$_0$ and H$_a$ for a $t$-test for the slope of a regression
          model in context.
    \item I can verify the LINER conditions for a $t$-test for slope given context
          information.
\end{itemize}
\end{learningtargetbox}

\begin{tocbox}
\renewcommand{\arraystretch}{1.5}
\begin{tabularx}{\linewidth}{clXc}
    \textbf{\#} & \textbf{Component} & \textbf{Description} & \textbf{Score} \\
    \hline
    1 & Warm-Up      & Spiral review: CI for slope, hypotheses preview & \hfill/6 \\
    2 & Guided Notes & $t$-test for slope, hypotheses, LINER conditions & \hfill/12 \\
    3 & Activity     & Setting Up a Significance Test for Slope & \hfill/15 \\
    4 & Exit Ticket  & Independent: procedure, hypotheses, conditions & \hfill/6 \\
    5 & Homework     & Practice setting up $t$-tests for slope & \hfill/12 \\
    \hline
       & \multicolumn{2}{l}{\textbf{Total}} & \hfill/51 \\
\end{tabularx}
\end{tocbox}

\end{document}
